% !PRINTED 53
\weberpage{53}
%<<CONT

\gesperrt{The $\mu$th derivative of the $\nu$th derivative is the $(\nu+\mu)$th derivative of the original function.}

Thus all the higher derivatives are obtained by forming from each preceding one, according to the rule (3), the first derivative:
\begin{align}
 f(x)&=a_0x^n+a_1x^{n-1}+a_2x^{n-2}+a_3x^{n-3}+\cdots,\notag\\
 f'(x)&=na_0x^{n-1}+(n-1)a_1x^{n-2}+(n-2)a_2x^{n-3}+\cdots, \tag{9}\\
 f''(x)&=n(n-1)a_0x^{n-2}+(n-1)(n-2)a_1x^{n-3}+\cdots. \notag
\end{align}

These derivatives assume a somewhat simpler form if one uses another notation, which is often in use and for certain purposes almost indispensable, and which we shall discuss in connection with this.

Because of the indeterminacy of the coefficients $a_0,a_1,\ldots,a_n$, it is clearly no restriction if we write a whole rational function of degree $n$ in the form
\begin{equation}
 f(x)=a_0x^n+B_1^{(n)}a_1x^{n-1}+B_2^{(n)}a_2x^{n-2}+\cdots+a_n,
\tag{10}
\end{equation}
or explicitly
\begin{equation}
 f(x)=a_0x^n+na_1x^{n-1}+{n(n-1)\over1\cdot2}a_2x^{n-2}+\cdots.
\tag{11}
\end{equation}
When a function $f(x)$ is represented in this way, we shall say that it is ``\gesperrt{written with the binomial coefficients}.''

A greater agreement is already shown by the formulae (9), which then read:
\begin{align}
 f(x)&=a_0x^n+na_1x^{n-1}+{n(n-1)\over1\cdot2}a_2x^{n-2}+\cdots,\notag\\
 {1\over n}f'(x)&=a_0x^{n-1}+(n-1)a_1x^{n-2}+{(n-1)(n-2)\over1\cdot2}a_2x^{n-3}+\cdots, \tag{12}\\
 {1\over n(n-1)}f''(x)&=a_0x^{n-2}+(n-2)a_1x^{n-3}+{(n-2)(n-3)\over1\cdot2}a_2x^{n-4}+\cdots,\notag
\end{align}
wherein the right-hand sides all likewise appear written with the binomial coefficients. We shall later become more acquainted with the usefulness of this notation, but already here we must emphasize that the choice of the one or the other
%CONT>>
